% META: {"id": "1SPE-DERGLOBAL-ALG-015", "chapitre": "1SPE-DERIVATION-GLOBAL", "type_objet": "algorithme", "capacites_codes": ["C3", "C4"], "programme_alignment": "OFFICIAL_ALGORITHM_EXAMPLE", "mode_creation": "ex_nihilo", "sources_inspiration": [], "status": "generated"}

\section{Algorithmique : la méthode de Newton}

% ============================================================
% STRATE 1 --- Probleme, modele, algorithme
% ============================================================

\margeAppui{
Le programme cite « Methode de Newton, en se limitant a des cas favorables »
parmi ses \textbf{exemples d'algorithme}. Il la cite, il ne l'impose pas.
}

\textbf{Le probleme.} L'equation $x^2 = 2$ n'a pas de solution decimale. On
cherche a en approcher la solution positive aussi finement qu'on veut, sans
disposer d'une formule.

\propriete[L'idee : remplacer la courbe par sa tangente]{
Soit $f$ derivable et $x_n$ une valeur approchee d'une solution de $f(x)=0$.
La tangente a $\mathcal{C}_f$ au point d'abscisse $x_n$ a pour equation
\[
  y = f'(x_n)(x - x_n) + f(x_n).
\]
Si $f'(x_n) \neq 0$, cette droite coupe l'axe des abscisses en un unique point,
que l'on prend pour valeur approchee suivante :
\[
  x_{n+1} = x_n - \frac{f(x_n)}{f'(x_n)}.
\]
}

\margeAppui{
\textbf{Cas favorable.} Ici $f(x) = x^2 - 2$, avec $f'(x) = 2x$. Sur
$[1\,;\,2]$, $f$ est strictement croissante ($f' > 0$) et sa courbe est au
dessus de chacune de ses tangentes : partant de $x_0 = 2$, la suite reste dans
$[\sqrt{2}\,;\,2]$ et decroit. C'est ce que « cas favorable » veut dire.
}

Pour $f(x) = x^2 - 2$, la formule se simplifie :
\[
  x_{n+1} = x_n - \frac{x_n^2 - 2}{2x_n} = \frac{1}{2}\left(x_n + \frac{2}{x_n}\right).
\]

\begin{python}
def newton(f, fprime, depart, iterations):
    """Valeurs approchees successives d'une solution de f(x) = 0."""
    if iterations < 0:
        raise ValueError("le nombre d'iterations doit etre positif")
    x = depart
    valeurs = [x]
    for _ in range(iterations):
        pente = fprime(x)
        if pente == 0:
            raise ZeroDivisionError("tangente horizontale : la methode s'arrete")
        x = x - f(x) / pente
        valeurs.append(x)
    return valeurs
\end{python}

\exemple{
Avec $f(x) = x^2 - 2$, $f'(x) = 2x$ et $x_0 = 2$, l'algorithme produit
\[
  x_0 = 2,\quad
  x_1 = \frac{3}{2},\quad
  x_2 = \frac{17}{12},\quad
  x_3 = \frac{577}{408},\quad
  x_4 = \frac{665\,857}{470\,832},
\]
soit en valeurs decimales
\[
  2\;;\; 1{,}5\;;\; 1{,}416\,666\ldots\;;\; 1{,}414\,215\,686\ldots\;;\;
  1{,}414\,213\,562\,374\ldots
\]
Or $\sqrt{2} = 1{,}414\,213\,562\,373\ldots$ : quatre iterations donnent douze
decimales exactes. Le nombre de decimales correctes double a peu pres a chaque
etape.
}

\propriete[Travail demande]{%
\begin{enumerate}
  \item Verifier a la main que $x_1 = \frac{1}{2}\left(2 + \frac{2}{2}\right)
        = \frac{3}{2}$.
  \item Compter, pour chaque $x_n$, le nombre de decimales communes avec
        $\sqrt{2}$.
  \item Reprendre avec $f(x) = x^2 - 5$ et $x_0 = 3$ : combien d'iterations
        faut-il pour obtenir dix decimales de $\sqrt{5}$ ?
  \item Que se passe-t-il si l'on part de $x_0 = 0$ ? Expliquer avec la
        tangente.
\end{enumerate}
}

\exemple{
\textbf{Reponses.} 2. Les nombres de decimales exactes sont $0$, $1$, $2$, $5$
puis $11$. 3. Quatre iterations suffisent depuis $x_0 = 3$ ; trois n'en
donnent que six. 4. En $x_0 = 0$ la
tangente est horizontale : elle ne coupe pas l'axe des abscisses, et
l'algorithme s'arrete -- c'est le cas defavorable que la formule signale par
une division par $f'(x_n) = 0$.
}

% BEGIN-VERIFY
% from fractions import Fraction as F
% from math import sqrt
%
% def newton(f, fprime, depart, iterations):
%     if iterations < 0:
%         raise ValueError("le nombre d'iterations doit etre positif")
%     x = depart
%     valeurs = [x]
%     for _ in range(iterations):
%         pente = fprime(x)
%         if pente == 0:
%             raise ZeroDivisionError("tangente horizontale : la methode s'arrete")
%         x = x - f(x) / pente
%         valeurs.append(x)
%     return valeurs
%
% f = lambda x: x*x - 2
% fp = lambda x: 2*x
% suite = newton(f, fp, F(2), 4)
% assert suite == [F(2), F(3, 2), F(17, 12), F(577, 408), F(665857, 470832)]
% # La formule simplifiee donne la meme suite.
% assert all(b == (a + 2/a)/2 for a, b in zip(suite, suite[1:]))
% # Cas favorable : la suite decroit et reste au-dessus de racine de 2.
% assert all(float(x) >= sqrt(2) for x in suite)
% assert all(float(a) > float(b) for a, b in zip(suite, suite[1:]))
% # Nombre de decimales exactes : 0, 1, 2, 5, 11.
% def decimales_exactes(x):
%     ecart = abs(float(x) - sqrt(2))
%     n = 0
%     while ecart < 10.0**(-(n + 1)):
%         n += 1
%     return n
% assert [decimales_exactes(x) for x in suite] == [0, 1, 2, 5, 11]
% # Racine de 5 depuis 3 : dix decimales en quatre iterations.
% cinq = newton(lambda x: x*x - 5, lambda x: 2*x, F(3), 4)
% assert abs(float(cinq[4]) - sqrt(5)) < 1e-10
% assert abs(float(cinq[3]) - sqrt(5)) > 1e-10
% # Depart en 0 : tangente horizontale, la methode s'arrete.
% try:
%     newton(f, fp, F(0), 1)
% except ZeroDivisionError:
%     pass
% else:
%     raise AssertionError("la tangente horizontale aurait du arreter la methode")
% END-VERIFY

% ============================================================
% STRATE 2 --- Erreurs frequentes
% ============================================================

\erreurFrequente{
\textbf{Croire que la methode converge toujours.} Elle ne converge que dans des
cas favorables. Une tangente horizontale ($f'(x_n) = 0$) l'arrete net ; un
mauvais point de depart peut l'envoyer vers une autre solution, ou la faire
osciller.
}

\erreurFrequente{
\textbf{Confondre valeur approchee et valeur exacte.} $\frac{577}{408}$ n'est
pas $\sqrt{2}$ : c'est un rationnel qui en donne six decimales. La methode
approche, elle ne demontre pas l'irrationalite.
}
